\chapter{特征工程与特征选择：什么对模型真正重要}

\section{特征工程的重要性}

特征太少可能无法描述问题，特征太多则可能引入噪声、增加过拟合风险并降低可解释性。特征工程需要在信息充分与结构简洁之间取得平衡。

\section{将领域知识编码为特征}

\begin{itemize}
  \item 时间特征：小时、星期、节假日、早晚高峰、寒暑假和特殊事件。
  \item 空间特征：道路等级、距中央商务区距离、兴趣点密度、站点可达性和路网中心性。
  \item 行为特征：司机历史接单率、乘客取消率、换乘等待时间和订单响应时间。
  \item 交互特征：雨天与高峰期、学校周边与放学时间、机场区域与夜间、补贴与接驾距离。
\end{itemize}

\section{用AI辅助特征重要性分析}

可以让 AI 使用 SHAP、基于树的特征重要性和相关系数等方法排列特征重要性。但需要特别关注领域知识与统计结论之间的冲突：统计上“不重要”的特征，在实际业务中可能不可缺少。

\section{实战练习：交通事故严重程度预测中的特征选择}

比较多种特征重要性方法的结果，结合交通领域知识决定最终保留哪些特征，并记录每项选择的理由。
